\documentclass[solutions,answers]{exam}

\usepackage{amsmath}
\usepackage{amssymb}
\usepackage{tabularx}
\usepackage{graphicx}
\usepackage{enumerate}
\usepackage{hhline}
\usepackage{tikz}


%######################################
%Various shortcuts
%Some may need package amssymb

%Shortcuts for various sets
\providecommand{\A}{\ensuremath{\mathfrak{A}}}%ideal https://www.overleaf.com/project/5ba957d84dc5ce49f189bc04A
\providecommand{\B}{\ensuremath{\mathcal{B}}}%basis
\providecommand{\C}{\ensuremath{\mathbb{C}}}%complex
\providecommand{\F}{\ensuremath{\mathbb{F}}}%finite field
\providecommand{\N}{\ensuremath{\mathbb{N}}}%natural
\providecommand{\RI}{\ensuremath{\mathcal{O}}}%ring of integers
\providecommand{\p}{\ensuremath{\mathfrak{P}}}%prime ideal
\providecommand{\Q}{\ensuremath{\mathbb{Q}}}%rationals
\providecommand{\R}{\ensuremath{\mathbb{R}}}%reals
\providecommand{\T}{\ensuremath{\mathcal{T}}}%topology
\providecommand{\Z}{\ensuremath{\mathbb{Z}}}%integers
\providecommand{\Zpos}{\ensuremath{\mathbb{Z}^{+}}}%positive integers


%Miscellaneous shortcuts
\providecommand{\lcm}{\ensuremath{\text{lcm}}}%least common multiple
\providecommand{\st}{\ensuremath{\backepsilon}}
\providecommand{\ud}{\ensuremath{\,\,\mathrm{d}}}%For derivatives
%######################################

%######################################
% Various operators
% some may need package amsmath

%Algebra
\providecommand{\amod}{\ensuremath{\diagup}}%Algebraic mod
\providecommand{\embed}{\ensuremath{\hookrightarrow}}%inclusion map
\providecommand{\gen}[1]{\ensuremath{\left<#1\right>}}%Group generated by #1
\providecommand{\idx}[2]{\ensuremath{\left[#1:#2\right]}}%index of #2 in #1
\providecommand{\im}{\ensuremath{\mathrm{im}}}%Image
\providecommand{\iso}{\ensuremath{\simeq}}%isomorphism
\providecommand{\normal}{\ensuremath{\vartriangleleft}}%Normal subgroup
\providecommand{\subgp}{\ensuremath{\le}}%Sub group
\providecommand{\vect}[1]{\ensuremath{\left\langle#1\right\rangle}}%Vector

%Logic
\providecommand{\land}{\ensuremath{\wedge}}
\DeclareMathSymbol\lnot{\mathbin}{symbols}{"3A}%p528
\providecommand{\lor}{\ensuremath{\vee}}
\providecommand{\lxor}{\ensuremath{\oplus}}

%Number Theory
\providecommand{\jacobi}[2]{\ensuremath{\left(\frac{#1}{#2}\right)}}
\providecommand{\legendre}[2]{\ensuremath{\left(\frac{#1}{#2}\right)}}

%Set Theory
\providecommand{\intersect}{\ensuremath{\bigcap}}
\providecommand{\less}{\ensuremath{\diagdown}}
\providecommand{\union}{\ensuremath{\bigcup}}

%Miscellaneous
\providecommand{\abs}[1]{\ensuremath{\left\lvert#1\right\rvert}}
\providecommand{\card}[1]{\ensuremath{\left\lvert#1\right\rvert}}
\providecommand{\ceiling}[1]{\ensuremath{\left\lceil#1\right\rceil}}
\providecommand{\define}{\ensuremath{\stackrel{\text{\tiny def}}{=}}}
\providecommand{\floor}[1]{\ensuremath{\left\lfloor#1\right\rfloor}}
\providecommand{\norm}[1]{\ensuremath{\left\lVert#1\right\rVert}}%Analysis norm
\providecommand{\order}[1]{\ensuremath{\left\lvert#1\right\rvert}}

%Asymptotic
\providecommand{\Oh}{\ensuremath{\mathcal{O}}}
\providecommand{\oh}{\ensuremath{o}}

%Environments




\header{COMP 2300: Function Practice}{}{Fall 2025}
\lfoot{}
\cfoot{}
\rfoot{\thepage}
\begin{document}

Let $f:\mathbb{Z}\to \mathbb{N}$. Be defined by $f(x)=|2x|+3$

\begin{enumerate}
\vfill
\item What is the Domain?\\
\vfill
\item What is the Range?\\
\vfill
\item What is the Codomain?\\
\vfill
\item Does $f$ have an inverse?\\
\vfill


\noindent Give examples of a function from $\mathbb{N} \times \mathbb{N}$ to $\mathbb{N}$ which is.

\item one to one, but not onto
\vfill
\item onto but not one to one
\vfill
\item both one to one and onto
\vfill
\item neither one to one nor onto
\vfill
\newpage

For all $i \in \mathbb{N}$, \\
Let $A_i=\{0,1,2,3,\ldots, i-1\}$
For example, $A_3=\{0,1,2\}$


\item Define an injective function from  $A_{10}$ to $A_{10}\setminus A_2$ or state that no such function exists.

\vfill
\item Define a surjective function from $A_5$ to $A_6\setminus A_2$ or state that no such function exists
\vfill
\item Define a function that is surjective and not injective from $A_8\to A_{10}\setminus A_4$ or state that one does not exist.
\vfill
\item Give an example of a function from $\mathbb{Z}$ to $\mathbb{N}$ that is onto but not injective.
\vfill
\item Determine if the function  $f:\mathbb{Z}\to \mathbb{Z}$ is injective, surjective, a bijection, and some of the above. $f(x)=|x+17|$.
\vfill
\item Let $f$ be a function from $\mathbb{Z}$ to $\mathbb{R}$ defined by $f(x)=3x^2$. Determine if the function  $f$ is injective, surjective, a bijection, and some of the above
\vfill
\end{enumerate}

\end{document}